\section{提案システムの設計}
本システムではバックエンドとフロントエンドに機能を分けて設計しており，帰宅推奨時刻の算出，画面表示，通知，投票集計といった処理が連携して動作する．
バックエンドでは，データベースに算出された帰宅推奨時刻の保存，Slackへの通知とメッセージの受信，投票結果の集計など，データの管理と外部連携を担当する．
フロントエンドは，画面表示および音声通知を担当するとともに，GitHub Actionsを用いてVercelにデプロイされたAPI Routeを定期実行し，帰宅推奨時刻の算出やSlack通知のロジックを呼び出す．

具体的なシステムフローは以下の通りである．
\begin{enumerate}
  \item GitHub ActionsでAPI Routeを定期実行し，帰宅推奨時刻を算出
  \item 有効な時刻が算出された場合，バックエンドのWeb APIを呼び出して保存・最多区間内のメンバに通知
  \item 画面に投票の選択肢を提示し，音声通知
  \item バックエンドがSlackイベントを通じて投票を受信しデータベースに保存
  \item 投票締切後，集計結果を画面表示に反映
\end{enumerate}

本システムでは，ユーザは提示される帰宅時刻の選択肢に投票するかどうか，最多投票の時刻に合わせて帰宅するかどうかを自由に決定できる設計としている．
投票したかどうかは他者に公開されず，最多投票の時刻に投票していたとしてもその時刻に合わせた帰宅を強制されない．
このため，コミュニティ内に行動を共にしたくない人がいる場合に理由を後付けして断りやすく，心理的負担を抑えられる．
さらに，画面提示は個別端末でも確認できるが全体に向けて提示され，予測帰宅時刻が近いとされる通知対象メンバも非公開とする．
この設計により，特定の人にのみ情報が知られる状況はなく，帰宅しない選択をしても不自然にならないよう配慮した．

% 
さらに，ユーザが負荷なく他端末から画面表示を確認できるように，アクセス可能な状態で提示される設計としている。
これにより，提示画面から遠い位置にいても状況を把握できるので，ユーザの負担を軽減できる．
一方で，個々の滞在者や各時刻への投票者，実際の帰宅時刻などの情報は公開せず，個人を特定できるデータが他者に見えないようにしている。
この設計により，ユーザはプライバシーを確保しつつ，必要な情報を負担なく確認できる環境が提供されている．